\section{Limitations}\label{sec:limitations}
% Rows: A1 (scope of the result), C5, C3, C6, C1, A7, A9, A3, B6; section E rows
% "our synthetic agent represents real agents" (DA QUALIFICARE) and "origin RPS = backend
% load in general" (DA QUALIFICARE). Never "benign" (section E: NOT SUPPORTED).

\paragraph{One testbed.}
All measurements come from one testbed: one virtual machine, one corpus, one cache
technology with LRU eviction, and one synthetic generator.\claim{A1}
We did not test other eviction policies or other hosts on the testbed (the simulator
explores other policies, \cref{sec:m-sim}), the primary design uses one cache size, and absolute values such as origin rates and miss ratios are properties of this
testbed.
All experiments use the cacheable chapter page; requests that no cache can absorb, such as
full-text search, are outside this study.
With the corrected exhaustive traversal each reachable-set size has one run of 5
repetitions per point, and
the smallest was measured on a different day from the other two; the earlier measurements
with the shared counter agree on the sign and the ordering.\claim{A1}

\paragraph{A synthetic agentic class.}
The generator matches the agentic clients of our honeypot in the width of the windowed
working set but not in its contiguity.\claim{C3}\claim{C5}
Measured the same way on both sides, 66.9\% of consecutive request pairs within a session
read adjacent chapters in the generator, against 4.25\% for honeypot agentic
clients.\claim{C5}
The unit of a session also differs: a k6 virtual user in the generator, a TCP connection
on the honeypot.\claim{C5}
The synthetic agentic class is therefore a controlled reachable set of a given width, not a
model of how agentic clients read a site.

\paragraph{One honeypot.}
The comparison with real clients rests on one public measurement site, one period and 31
clients of the agentic class with at least 5 requests.\claim{C3}
Clients are classed by their declared user agent, which we do not verify.
The honeypot results describe that population in that period and do not characterise
agentic traffic on the Web in general.\claim{C3}\claim{C6}

\paragraph{Origin requests as a measure of cost.}
The relation between origin requests and PostgreSQL CPU is established on this testbed
only.\claim{C1}
We do not claim that origin requests measure backend load in general.\claim{C1}
We also did not measure where origin load turns into human latency: the queue at each
stage was never measured (\cref{sec:m-human}).\claim{B6}

\paragraph{Secondary series and shared mapping.}
The series that varies the exhaustive volume comes from an earlier campaign with a
different mapping and window (\cref{app:a7}), so it is not paired with the primary
design.\claim{A7}

\paragraph{Runs that keep the shared traversal counter.}
The admission-policy frontier (\cref{sec:r-frontier}), the CPU validation of origin
requests, the measurement of the cache capacity, the window-length check and the secondary
exhaustive series (\cref{app:a7}) come from runs made before the exhaustive traversal was
corrected, and were not rerun.\claim{A6}\claim{C1}\claim{C2}\claim{C7}\claim{A7}
For the frontier, a chapter came up in the shared traversal every 91.6\,s (16,954 chapters
at 185 req/s), and the exhaustive class requested it at each pass with probability
$\alpha = 0.35$.\claim{A6}
Taking the characteristic time of each policy as the cache capacity over its origin rate,
$T \approx 5{,}263$ objects divided by the origin rate, gives 65\,s with no policy; 131,
96, 82, 77 and 70\,s for deferral budgets 1, 2, 3, 4 and 6; and 280 and 293\,s for the two
blocking policies.\claim{A6}
A return after 91.6\,s therefore falls within $T$ only under budgets 1 and 2, where
exhaustive hits may be inflated; with no policy and budgets 3 to 6, $T$ is shorter than the
return.\claim{A6}
Under the blocking policies an exhaustive miss receives an uncacheable error, so the
exhaustive class never hits its own earlier requests.\claim{A4}\claim{D2}
We did not measure how much budgets 1 and 2, and hence the shape of the frontier, are
affected.\claim{A6}
Except for the capacity runs, these runs also predate the recording of full configurations
on 21 September 2026.
In the exhaustive series, the only one where their rates enter a result, $\alpha$ is read
from the result files and $\lambda$ from the campaign script, and $\lambda$ is checked
against the request rate that k6 measured.\claim{A7}
All primary results at scope 0.02 and across the three reachable sets come from runs with
the corrected traversal and a recorded configuration.\claim{A1}\claim{A3}\claim{A9}
